\documentclass[12pt]{article}
\usepackage{amsmath}
\usepackage{amssymb}
\usepackage{geometry}
\usepackage{enumitem}
\geometry{a4paper, margin=1in}

\begin{document}

\begin{center}
    \textbf{University of Eswatini} \\
    \textbf{Department of Computer Science} \\
    \textbf{Examination (Main)} \\
    2025/2026 \\
    \textbf{FIRST SEMESTER} \\
    \vspace{1cm}
    \textbf{INTRODUCTION TO COMPUTER SCIENCE} \\
    \textbf{Course Code: CSC111} \\
    \vspace{1cm}
    \textbf{Time Allowed: Three (3) Hours}
\end{center}

\noindent
\textbf{Instructions:}
\begin{enumerate}
    \item Answer \textbf{ALL} questions in \textbf{Section A} (50 marks).
    \item Answer \textbf{EITHER} all questions in \textbf{Section B} (50 marks) \textbf{OR} all questions in \textbf{Section C} (50 marks).
    \item Do not write anything on the question paper.
    \item You are not allowed to open this paper until instructed by the invigilator.
    \item Any type of calculator is not allowed.
\end{enumerate}

\vspace{1cm}
\newpage
\section*{SECTION A: COMPULSORY (50 marks)}
\textit{Answer all questions in this section.}

\begin{enumerate}[left=0pt, label=\bfseries Q\arabic*.]
    \item \textbf{Computer Science Fundamentals} \hfill (10 marks)\\
          List and explain the five distinct computing disciplines identified by the Association for Computing Machinery (ACM).

    \item \textbf{Computer Generations} \hfill (10 marks)\\
          Describe the five generations of computers, highlighting the key technology and characteristics of each generation.

    \item \textbf{System Development Life Cycle (SDLC)} \hfill (10 marks)\\
          List and explain the six phases of the System Development Life Cycle (SDLC).

    \item \textbf{Computer Literacy and Professionals} \hfill (10 marks)\\
          Differentiate between the following:
          \begin{enumerate}[label=(\alph*)]
              \item Computer literacy and computer competency
              \item Computer professional and computer user (give two examples of each)
          \end{enumerate}

    \item \textbf{Binary Arithmetic and Storage} \hfill (10 marks)\\
          \begin{enumerate}[label=(\alph*)]
              \item Convert the binary number 101101 to decimal. \hfill (3 marks)
              \item Convert the hexadecimal number A3F to binary. \hfill (3 marks)
              \item Convert 2.5 GB to MB. \hfill (4 marks)
          \end{enumerate}
\end{enumerate}

\newpage

\section*{SECTION B (50 marks)}
\textit{Answer ALL questions in this section.}

\begin{enumerate}[left=0pt, label=\bfseries Q\arabic*.]
    \item \textbf{Algorithms and Flowcharts} \hfill (15 marks)\\
          \begin{enumerate}[label=(\alph*)]
              \item Define an algorithm and list atleast four characteristics of a good algorithm. \hfill (5 marks)
              \item Write an algorithm to calculate the area and circumference of a circle given the radius. \hfill (5 marks)
              \item Draw a flowchart to represent the algorithm for finding the largest of three numbers. \hfill (5 marks)
          \end{enumerate}

    \item \textbf{Python Programming I} \hfill (15 marks)\\
          \begin{enumerate}[label=(\alph*)]
              \item Write a Python program to calculate the area of a triangle. The program should ask the user for the \texttt{base} and \texttt{height} and then display the \texttt{area}. \hfill (5 marks)
              \item Write a Python program that takes a string from the user and does the following:
                    \begin{itemize}
                        \item Prints the string in uppercase.
                        \item Prints the string in lowercase.
                        \item Prints the length of the string.
                        \item Prints the string with the first and last characters removed.
                    \end{itemize}
                    \hfill (10 marks)
          \end{enumerate}

    \item \textbf{Systems Analysis and Design} \hfill (20 marks)\\
          \begin{enumerate}[label=(\alph*)]
              \item Define the term "system" and list its key elements. \hfill (5 marks)
              \item Explain why systems fail. List any five reasons. \hfill (5 marks)
              \item What is the importance of documentation in system development? \hfill (5 marks)
              \item Differentiate between system software and application software. \hfill (5 marks)
          \end{enumerate}
\end{enumerate}

\newpage

\section*{SECTION C (50 marks)}
\textit{Answer ALL questions in this section.}

\begin{enumerate}[left=0pt, label=\bfseries Q\arabic*.]
    \item \textbf{8-bit Binary Arithmetic} \hfill (15 marks)\\
          Assume an 8-bit binary machine where the \( 8^{th} \) bit is the sign bit (0 for positive, 1 for negative). Only 7 bits are available for magnitude. Given: \( A = 20_{10}, \, B = 2B_{16}, \, C = 36_{8}, \, D = -12_{10} \), answer the following:
          \begin{enumerate}[label=(\alph*)]
              \item Express \( D \) as an 8-bit number using 2’s complement representation. \hfill (3 marks)
              \item Calculate the binary sum of A and B. Provide the result in 8-bit binary. \hfill (4 marks)
              \item Calculate \( A - B \) using the 2’s complement method. Provide the result in 8-bit binary. \hfill (4 marks)
              \item Can the number \( -132_{10} \) be represented in this 8-bit machine? Justify your answer. \hfill (4 marks)
          \end{enumerate}

    \item \textbf{Python Programming II} \hfill (15 marks)\\
          \begin{enumerate}[label=(\alph*)]
              \item Write a Python program that acts as a simple calculator. It should ask the user for two numbers and an operation (+, -, *, /) and then perform the operation and display the result. \hfill (8 marks)
              \item Write a Python program to generate the multiplication table of 8 up to 12 times (i.e., \(8 \times 1 = 8\) to \(8 \times 12 = 96\)). \hfill (7 marks)
          \end{enumerate}

    \item \textbf{Computer Architecture and Number Systems} \hfill (20 marks)\\
          \begin{enumerate}[label=(\alph*)]
              \item Explain the stored program concept in the Von Neumann architecture. \hfill (5 marks)
              \item Differentiate between RAM and ROM. \hfill (5 marks)
              \item What is the role of the control unit in a CPU? \hfill (5 marks)
              \item Convert the octal number 735 to binary. \hfill (5 marks)
          \end{enumerate}
\end{enumerate}

\end{document}